\section{Pattern Recognition}

Sau khi phân rã bài toán, bước tiếp theo của tư duy tính toán là tìm ra
những \textbf{pattern} lặp đi lặp lại giữa các bài toán con. Điều này có
ý nghĩa quan trọng, vì nhiều khi các vấn đề trông có vẻ khác nhau ở
bề mặt nhưng thực chất lại có chung một cấu trúc xử lý bên dưới. Khi
nhận ra các quy luật đó, ta có thể tái sử dụng thiết kế, dữ liệu và
logic thay vì xây dựng mọi thứ theo cách rời rạc~\cite{spronck2023, dejesus2020}.

\textbf{Đầu tiên} là mẫu \textbf{context-aware decision making}. Dù là
feed gợi ý địa điểm, bubble map hay lựa chọn nhiệm vụ kế tiếp trong
storyline thì tất cả đều dựa trên cùng một câu hỏi nền tảng: với trạng
thái hiện tại của người dùng và môi trường xung quanh, hệ thống nên ưu
tiên điều gì trước. Điều này có nghĩa là nhiều tính năng khác nhau thực
ra cùng chia sẻ một lớp tư duy ra quyết định theo ngữ cảnh. Từ pattern
này, nhóm đi tới quyết định rằng feed và bubble map không được phép dùng
hai công thức khác nhau, mà phải dùng cùng một lõi suitability score để
tránh việc hai bề mặt hiển thị đưa ra hai thông điệp mâu thuẫn. Nếu
không nhận ra pattern này từ sớm, rất dễ xảy ra tình huống mỗi màn hình
được xây theo một logic riêng: feed xếp theo mức hấp dẫn nội dung, bản
đồ phóng to theo khoảng cách, còn storyline lại ưu tiên theo cảm hứng
riêng của AI. Khi đó người dùng sẽ không còn thấy một hệ thống thống
nhất, mà chỉ thấy nhiều tính năng đứng cạnh nhau nhưng không nói cùng
một ngôn ngữ quyết định.

\textbf{Kế đến} là mẫu \textbf{retrieve → explain → present}. Đây chính là một pattern
phổ biến trong xây dựng hệ thống AI đó chính Retrieval-Augmented Generation (RAG) sau đó giải
thích giá trị của địa điểm đó và cuối cùng trình bày nó thành card, feed
hoặc bong bóng bản đồ. Trong storyline, hệ thống cũng truy xuất bối cảnh
văn hóa, sinh ra mô tả nhiệm vụ và cuối cùng hiển thị nó cho người dùng
dưới dạng các bước hành động. Điều này cho thấy dù giao diện đầu ra khác
nhau, cấu trúc xử lý thông tin vẫn rất tương đồng. Từ đó mà nhóm thiết
kế AI theo hướng lấy dữ liệu văn hóa có nguồn gốc trước, sau đó mới sinh
phần giải thích hoặc nhiệm vụ bằng mô hình ngôn ngữ~\cite{bevietnam_spec}.
Điều này có ý nghĩa thực tế rất lớn, vì nó ngăn hệ thống rơi vào cách
làm phổ biến nhưng rất yếu là đưa prompt chung chung cho mô hình rồi hy
vọng mô hình ``sẽ tự biết'' nên nói gì. Thay vào đó, từng khối tính năng
đều phải đi qua cùng một cấu trúc: lấy đúng dữ liệu nền, sinh nội dung
dựa trên dữ liệu đó, rồi mới trình bày ra giao diện theo hình thức phù
hợp.

\textbf{Cuối cùng} là mẫu \textbf{verify → update state}. Sau mỗi hành động
quan trọng của người dùng, ví dụ như hoàn thành một nhiệm vụ hay gửi một
bức ảnh minh chứng, hệ thống đều phải kiểm tra dữ liệu đầu vào, xác minh
tính hợp lệ rồi mới cập nhật trạng thái tiến trình. Mẫu này lặp lại rất
nhiều trong backend, AI service và giao diện người dùng, từ đăng nhập,
gửi capture cho tới hoàn thành task. Pattern này dẫn nhóm tới yêu cầu là
mọi hành động cập nhật tiến trình phải đi qua lớp validation và không để
AI ghi trực tiếp vào product state. Đây là quyết định quan trọng ở
giai đoạn pilot, vì nếu bỏ qua lớp kiểm tra này thì các tính năng liên
quan đến tài khoản, tiến trình storyline và capture sẽ rất dễ rơi vào
tình trạng không nhất quán, đặc biệt khi hệ thống còn đang dùng nhiều
nguồn dữ liệu và nhiều bề mặt client khác nhau.

Việc nhận diện các mẫu lặp này có hậu quả kỹ thuật rất rõ ràng. Trước hết, nó giúp nhóm chuẩn hóa cách thiết kế API, vì nhiều luồng có thể
dùng chung cấu trúc request/response và cách quản lý trạng thái. Kế đó, nó giúp nhóm chuẩn hóa cách viết service và repository, bởi dù dữ
liệu thuộc về \texttt{place}, \texttt{feed} hay \texttt{task}, ta vẫn
có thể tổ chức theo các tầng rất giống nhau. Sau cùng, nó tạo tiền đề để
AI không hoạt động như một hộp đen độc lập, mà chỉ là một mắt xích trong
một pipeline có pattern rõ ràng: nhận dữ liệu đầu vào, suy luận trong
phạm vi kiểm soát và trả về output có cấu trúc.
